% ============================================================================
%  Actividad 4 · Introduccion y nota de herramienta
%
%  Propiedad de US-UX-07. Cubre el apartado 2 de la rubrica y el criterio 7 de
%  la historia: la nota que justifica construir en el stack del producto,
%  anclada a las cuatro propiedades que la rubrica pide de un prototipo de alta
%  fidelidad y al objetivo de aprendizaje OEA 2.2.
% ============================================================================

% ============================================================================
\uxsection{Introducción}

Este documento presenta el \textbf{prototipo de alta fidelidad} de \textbf{Karisma Data}, el Portal
Centralizado de Datos Financieros, y la \textbf{guía de estilos} del sistema de diseño con el que
está construido. Es la cuarta entrega de una investigación acumulativa y la primera en la que el
proyecto deja de describir un producto para mostrarlo funcionando.

\textbf{Continuidad de la serie.} La Actividad 1 delimitó el problema, la audiencia institucional y
las ocho personas que representan los perfiles de uso. La Actividad 2 convirtió esas personas en
escenarios y mapas de recorrido, y de ellos derivó cinco principios rectores. La Actividad 3 situó
al portal frente a sus competidores, contrastó la estructura propuesta con el modelo mental de quien
la usará mediante un ejercicio de card sorting de treinta y cinco tarjetas, y tradujo ese resultado
en una arquitectura de información de cuatro categorías y dieciséis ramas. Esta actividad toma ese
mapa como contrato y lo convierte en interfaz: cada ruta del prototipo se ancla a una rama de aquel
mapa, y la correspondencia se publica completa.

\textbf{Qué se entrega.} Siete pantallas de alta fidelidad, construidas en el entorno del producto y
recorribles de principio a fin; la correspondencia entre las dieciséis ramas del mapa y las rutas
del portal; los cuatro estados no felices declarados pantalla por pantalla; una guía de estilos con
sus secciones puntuadas y sus dos temas verificados bajo el mismo listón de contraste en modo claro
y en modo oscuro; cinco flujos de tarea completos en diseño de alta fidelidad; la matriz de alcance
que declara hasta dónde llega cada pantalla; y una validación sobre capturas reales del prototipo,
con sus tres ejes de refinamiento y la medición de cada uno.

\textbf{Sobre la validación.} Los ocho evaluadores que aparecen en la sección de validación
son \textbf{evaluadores prototipo}: perfiles sintéticos condicionados por las ocho personas de la
Actividad 1, en continuidad con el método que sostuvo el card sorting de la Actividad 3. No son
participantes humanos y sus respuestas no son datos de campo. La validación con personas reales
ocurre en la prueba de usabilidad de la Actividad 5, cuando el objeto de evaluación ya es una
interfaz operable y la medición resulta interpretable. Esta precisión abre la sección
correspondiente y se repite ahí, porque confundir las dos cosas equivaldría a fabricar un resultado.

% ============================================================================
\section{Sobre el medio del prototipo}

El equipo trabajó en dos superficies. Las pantallas se construyeron \textbf{en el entorno del
producto} --- Nuxt 4 sobre FastAPI, con PostgreSQL detrás --- y en un \textbf{archivo de diseño} se
resolvieron los cinco flujos de tarea, la identidad visual y el inventario de componentes con sus
variantes. La sección 1 explica qué demuestra cada una y por qué ninguna sustituye a la otra.

Construir en el entorno del producto merece justificarse, porque lo habitual en esta etapa es un
archivo de lienzos. La tabla contrasta las cuatro propiedades que se le piden a un prototipo de alta
fidelidad con lo que entrega una interfaz que se ejecuta.

\begin{uxtabla}{|L{5.4cm}|Y|}{Las cuatro propiedades de un prototipo de alta fidelidad frente a una interfaz operable}
  \uxheadrow \thd{Propiedad esperada} & \thd{Cómo la cumple una interfaz construida en el entorno del producto} \\
  \enquote{Ofrecer una representación lo más fiel al producto final} & La representación no se aproxima al producto final: es el producto final. Mismo motor de composición, misma tipografía cargada por el navegador, mismos datos sintéticos y mismo sistema de tokens \\
  \enquote{Permitir probar y evaluar interacciones y funcionalidades} & Las interacciones no se simulan entre lienzos: se ejecutan. Los filtros filtran, la serie de la gráfica hace acercamiento, la transmisión del asistente se cancela de verdad y la sesión cambia lo que cada rol ve \\
  \enquote{Facilitar la obtención de comentarios detallados de los usuarios acerca de la estética visual} & Se obtienen los mismos comentarios sobre estética y además sobre comportamiento, densidad de información y respuesta a la espera, que un lienzo estático no puede provocar \\
  \enquote{Ayudar a las personas interesadas a tener un mejor entendimiento del producto final} & No requiere que quien lo revisa imagine el resultado: se abre en un navegador y se recorre, con su sesión y sus permisos \\
\end{uxtabla}

Un prototipo que \emph{es} la interfaz la representa de la manera más completa disponible, y nada
se pierde por construirlo así. Este documento entrega el sistema de tokens
completo, el inventario de componentes con sus estados, la retícula y la matriz de contraste
verificada por cálculo, y añade lo que un lienzo no puede entregar: pantallas que responden. Cada
captura de la sección de prototipos es una imagen de la aplicación en funcionamiento, tomada a
resolución fija con un guion reproducible que se documenta en el anexo; cada lámina de flujo lleva
en su pie la palabra diseño, para que ninguna se lea como una captura de la aplicación.

% ============================================================================
\section{Cómo se lee este documento}

\begin{uxlist}
  \lead{Sección 1. Método de prototipado}{Con qué se construyó, qué significa \enquote{desplegado}
    en esta entrega y cómo se toman las capturas. Incluye la cadena que va del sistema de diseño a
    la página impresa y por qué corre en un solo sentido.}
  \lead{Sección 2. Los siete prototipos}{Las pantallas, una por una, con su figura. La
    correspondencia entre las dieciséis ramas del mapa de la Actividad 3 y las rutas del portal, y
    la matriz de los cuatro estados no felices.}
  \lead{Sección 3. Guía de estilos}{Identidad ---la marca del informe y la del portal, con su
    geometría medida y sus tres variantes---, tipografía, la especificación del sistema de tokens
    con sus dos temas y sus dos modos, componentes de interfaz e iconografía, cada uno con su
    propio subtítulo. Ahí viven las tres matrices del sistema: los cuarenta y cuatro pares de la
    paleta base, el contraste del tema institucional en los dos modos y la separación bajo las tres
    dicromacias.}
  \lead{Sección 4. Validación del prototipo}{El protocolo con los ocho evaluadores prototipo, las
    dos tareas heredadas de la prueba de árbol de la Actividad 3, y los tres ejes de refinamiento
    del sistema, cada uno con su regla, su línea base y su valor verificado.}
  \lead{Sección 5. Especificación de alcance}{Hasta dónde llega cada pantalla y cada historia de
    usuario, en cuatro estados, sin dejar ninguna sin declarar.}
  \lead{Sección 6. Los cinco flujos de tarea}{Cada flujo con su persona, su objetivo, su número de
    pantallas y su lámina, y la verificación técnica del archivo de diseño.}
  \lead{Sección 7. Cierre}{Qué se construyó en cifras y con qué mecanismo se comprueba cada pieza,
    qué cambió por evidencia y qué se lleva la Actividad 5.}
  \lead{Anexo}{Guion de capturas, inventario de figuras y trazabilidad entre los criterios de la
    historia de usuario, las secciones de este documento y su evidencia.}
\end{uxlist}

Se documentan siete pantallas y no las cinco que el mapa de la Actividad 3 marcaba como núcleo. Las
dos que se añaden ---exportación y asistente--- son las que sostienen el resto del recorrido: sin la
primera, encontrar un dato no termina en nada que el analista pueda llevarse; sin la segunda, la
promesa de preguntar en lenguaje natural se queda escrita y sin demostrar.
